\documentclass{article}
\usepackage{amsmath,amssymb,amsthm}
\usepackage{booktabs}
\usepackage{hyperref}
\usepackage{listings}
\usepackage{xcolor}
\usepackage{longtable}
\usepackage{array}
\usepackage[margin=1in]{geometry}

\newcommand{\hex}[1]{\texttt{\$#1}}
\newcommand{\opcode}[1]{\texttt{#1}}
\newcommand{\Beff}{B_{\mathrm{eff}}}
\newcommand{\Prep}{P_{\mathrm{rep}}}

% Soup class symbols (weight-based)
\newcommand{\wa}{w_a}
\newcommand{\wf}{w_f}
\newcommand{\wx}{w_x}
\newcommand{\wy}{w_y}
\newcommand{\wix}{w_{ix}}
\newcommand{\wiy}{w_{iy}}
\newcommand{\wdx}{w_{dx}}
\newcommand{\wdy}{w_{dy}}
\newcommand{\wh}{w_h}
\newcommand{\ws}{w_s}
\newcommand{\wu}{w_u}
\newcommand{\pmin}{p_{\min}}

\newtheorem{definition}{Definition}
\newtheorem{proposition}{Proposition}
\newtheorem{theorem}{Theorem}
\newtheorem{remark}{Remark}

\lstset{
  basicstyle=\ttfamily\small,
  columns=fullflexible,
  frame=single,
  backgroundcolor=\color{gray!10},
}

\title{Soup Class Analysis of the 6502 Replicator Mining Problem}
\author{Ian Holmes}
\date{April 2026}

\begin{document}
\maketitle

\begin{abstract}
We partition the 256 possible byte values into equivalence classes
(``soup classes'') based on their functional role in the six core
self-replicator loop variants of the 6502life cellular automaton.
We develop a probability model parameterized by per-class biases,
derive closed-form expressions for the replicator mining probability
as a function of insert count, and characterize the Pareto frontier
in the tradeoff space of mining time, insert diversity, and soup entropy.
\end{abstract}

\tableofcontents

%----------------------------------------------------------------------
\section{The Six Core Variants}
\label{sec:variants}

A minimal 6502life replicator consists of a \emph{copy loop body}
(6~core bytes) plus a \emph{branch opcode} and \emph{offset byte}
that close the loop.  The branch opcode is drawn from the set
$\{$\hex{10}, \hex{30}, \hex{50}, \hex{70}, \hex{90}, \hex{B0},
\hex{D0}, \hex{F0}$\}$ (the eight conditional branches:
BPL, BMI, BVC, BVS, BCC, BCS, BNE, BEQ).

The six core loop variants are:

\begin{center}
\begin{tabular}{@{}llll@{}}
\toprule
\textbf{Name} & \textbf{Family} & \textbf{Core bytes (hex)} & \textbf{Description} \\
\midrule
DEX     & X-std   & B5 00 9D 00 04 CA & \opcode{LDA \$00,X; STA \$0400,X; DEX} \\
INX     & X-std   & B5 00 9D 00 04 E8 & \opcode{LDA \$00,X; STA \$0400,X; INX} \\
DEY     & Y-std   & B7 00 99 00 04 88 & \opcode{LAX \$00,Y; STA \$0400,Y; DEY} \\
INY     & Y-std   & B7 00 99 00 04 C8 & \opcode{LAX \$00,Y; STA \$0400,Y; INY} \\
INX3FF  & X-shift & B5 00 E8 9D FF 03 & \opcode{LDA \$00,X; INX; STA \$03FF,X} \\
INY3FF  & Y-shift & B7 00 C8 99 FF 03 & \opcode{LAX \$00,Y; INY; STA \$03FF,Y} \\
\bottomrule
\end{tabular}
\end{center}

Each core can appear in any of 3 \emph{rotations} within the loop
(starting from each of the 3 instructions), and the loop can contain
$k \geq 0$ safe insert bytes between core instructions.

%----------------------------------------------------------------------
\section{Byte Classification}
\label{sec:classification}

\subsection{Core byte inventory}

We identify the \emph{distinct} byte values required by each variant,
annotated with the soup class to which each belongs (defined in
Section~\ref{sec:classes}):

\begin{center}
\begin{tabular}{@{}llll@{}}
\toprule
\textbf{Byte} & \textbf{Class} & \textbf{Required by} & \textbf{Multiplicity} \\
\midrule
\hex{00} & $\wa$ & All 6 variants & 2 (std), 1 (shift) \\
\hex{B5} & $\wx$ & DEX, INX, INX3FF & 1 \\
\hex{9D} & $\wx$ & DEX, INX, INX3FF & 1 \\
\hex{B7} & $\wy$ & DEY, INY, INY3FF & 1 \\
\hex{99} & $\wy$ & DEY, INY, INY3FF & 1 \\
\hex{04} & $\wf$ & DEX, INX, DEY, INY & 1 \\
\hex{E8} & $\wix$ & INX, INX3FF & 1 \\
\hex{C8} & $\wiy$ & INY, INY3FF & 1 \\
\hex{CA} & $\wdx$ & DEX only & 1 \\
\hex{88} & $\wdy$ & DEY only & 1 \\
\hex{03} & $\wh$ & INX3FF, INY3FF & 1 \\
\hex{FF} & $\wh$ & INX3FF, INY3FF & 1 \\
\bottomrule
\end{tabular}
\end{center}

\subsection{Variant--byte requirements}

Each variant's 6 core bytes, annotated by class:
\begin{center}
\begin{tabular}{@{}ll@{}}
\toprule
\textbf{Variant} & \textbf{Core bytes (class)} \\
\midrule
DEX     & B5($x$)\; 00($a$)\; 9D($x$)\; 00($a$)\; 04($f$)\; CA($dx$) \\
INX     & B5($x$)\; 00($a$)\; 9D($x$)\; 00($a$)\; 04($f$)\; E8($ix$) \\
DEY     & B7($y$)\; 00($a$)\; 99($y$)\; 00($a$)\; 04($f$)\; 88($dy$) \\
INY     & B7($y$)\; 00($a$)\; 99($y$)\; 00($a$)\; 04($f$)\; C8($iy$) \\
INX3FF  & B5($x$)\; 00($a$)\; E8($ix$)\; 9D($x$)\; FF($h$)\; 03($h$) \\
INY3FF  & B7($y$)\; 00($a$)\; C8($iy$)\; 99($y$)\; FF($h$)\; 03($h$) \\
\bottomrule
\end{tabular}
\end{center}

Each variant also requires one branch opcode (any of 8, all in class $a$)
and one offset byte (a specific value whose probability is $\pmin = \wu/W = 1/W$).

\subsection{Branch opcodes}

The 8 branch opcodes are:
\hex{10} (BPL), \hex{30} (BMI), \hex{50} (BVC), \hex{70} (BVS),
\hex{90} (BCC), \hex{B0} (BCS), \hex{D0} (BNE), \hex{F0} (BEQ).
These are grouped with \hex{00} into class~$\wa$ because
every variant requires exactly one branch opcode,
and \hex{00} appears in all 6 cores ---
both byte values belong to the set needed by \emph{all} variants.

The offset byte value is determined by the loop length
(e.g.\ \hex{F8} for an 8-byte loop, \hex{F7} for 9~bytes, etc.).
This value is almost certainly not in any elevated soup class, so
its probability equals the \emph{minimum} byte probability:
\[
  \pmin \;=\; \frac{\wu}{W}
  \;=\; \frac{1}{W},
\]
where $W = \sum_c n_c\,w_c$ is the total weight (Section~\ref{sec:model})
and we set $\wu = 1$ as the baseline.
This acts as a natural floor: concentrating the soup on core bytes
increases $W$, which \emph{decreases} $\pmin$ --- a self-regulating
mechanism.
\label{sec:offset}

\subsection{Safe insert bytes}
\label{sec:safe}

A \emph{safe insert} is a single-byte opcode that can execute between
core instructions without disrupting the copy loop.  We use a
\textbf{conservative} definition: the byte must be safe regardless of
rotation and regardless of branch opcode, meaning it must not modify
the accumulator~(A), the index register used by the variant (X~or~Y),
or any of the four condition flags (N, V, C, Z).

Under these constraints, the safe single-byte opcodes are those that
either perform no operation or only modify the stack pointer~S, the
interrupt-disable flag~I, or the decimal flag~D (none of which affect
any conditional branch).

\paragraph{Conservative 1-byte safe set (both families):}
\begin{itemize}
\item \hex{EA} --- \opcode{NOP} (official)
\item \hex{1A}, \hex{3A}, \hex{5A}, \hex{7A}, \hex{DA}, \hex{FA}
      --- undocumented 1-byte NOPs
\item \hex{48} --- \opcode{PHA} (push A; modifies only S)
\item \hex{08} --- \opcode{PHP} (push P; modifies only S)
\end{itemize}
That gives $n_s = 9$ bytes in the safe class.

\paragraph{Remark on \opcode{PHA}/\opcode{PHP}.}
Pushing onto the stack is harmless for the copy loop itself, but
repeated pushes without corresponding pulls will eventually overflow
the stack page (\hex{100}--\hex{1FF}).  Since the copy loop runs for
at most 256 iterations, the number of stack overflows is bounded, and
the stack wraps within its page.  The writes to the stack page are
within the cell's own memory and do not corrupt the source data
(zero page \hex{00}--\hex{FF}).  We therefore include \opcode{PHA}
and \opcode{PHP} in the safe set.

%----------------------------------------------------------------------
\section{Soup Classes: the Power Set Structure}
\label{sec:classes}

\subsection{Motivation}

The soup classes arise from the \textbf{power set} of the 6~core
variants.  Each byte value belongs to the \emph{subset of variants
that require it}.  Bytes needed by more variants are more valuable
to enrich because they contribute to more terms of the probability
sum.  For example, \hex{04} (class $\wf$, needed by 4 standard
variants) is more valuable per byte than \hex{CA} (class $\wdx$,
needed only by DEX).

\subsection{The 10 classes}

We partition all 256 byte values into 10 disjoint soup classes.
Each class $c$ has a per-byte weight $w_c$ and a size $n_c$.

\begin{definition}[Soup class partition]
\label{def:partition}
\begin{center}
\begin{tabular}{@{}clcll@{}}
\toprule
\textbf{Class} & \textbf{Symbol} & $n_c$ & \textbf{Bytes} & \textbf{Why} \\
\midrule
All + branches & $\wa$ & 9 & 00, 10, 30, 50, 70, 90, B0, D0, F0 &
  \hex{00} in all 6 cores; 8 branches for all \\
Standard page  & $\wf$ & 1 & 04 &
  Page byte for 4 std variants \\
X load/store   & $\wx$ & 2 & B5, 9D &
  \opcode{LDA zp,X}, \opcode{STA abs,X} \\
Y load/store   & $\wy$ & 2 & B7, 99 &
  \opcode{LAX zp,Y}, \opcode{STA abs,Y} \\
X increment    & $\wix$ & 1 & E8 &
  \opcode{INX} --- INX and INX3FF \\
Y increment    & $\wiy$ & 1 & C8 &
  \opcode{INY} --- INY and INY3FF \\
X decrement    & $\wdx$ & 1 & CA &
  \opcode{DEX} --- DEX variant only \\
Y decrement    & $\wdy$ & 1 & 88 &
  \opcode{DEY} --- DEY variant only \\
Shifted address & $\wh$ & 2 & 03, FF &
  Page \hex{03} + addr \hex{FF} for shifted variants \\
Safe inserts   & $\ws$ & $n_s$ & NOPs, PHA, PHP, etc. &
  Harmless when executed \\
Unsafe         & $\wu{=}1$ & $n_u$ & everything else &
  Background/fatal, fixed at weight 1 \\
\bottomrule
\end{tabular}
\end{center}
\end{definition}

With $n_s = 9$ (conservative safe set):
$n_u = 256 - 9 - 1 - 2 - 2 - 1 - 1 - 1 - 1 - 2 - 9 = 227$.

Verification: $9 + 1 + 2 + 2 + 1 + 1 + 1 + 1 + 2 + 9 + 227 = 256$.\ \checkmark

\begin{remark}
The key difference from a naive partition is that \hex{00} and the
8 branch opcodes share class~$\wa$.  This is correct because \hex{00}
appears in every core variant and every branch opcode is interchangeable,
so they have identical enrichment value.
\end{remark}

%----------------------------------------------------------------------
\section{Probability Model}
\label{sec:model}

\subsection{Parameterization}

Each byte $b \in \{0,\ldots,255\}$ is drawn independently from a
biased soup distribution parameterized by unnormalized per-byte
\emph{weights}.  All bytes in class $c$ share the same weight $w_c$.
The unsafe class is fixed at $\wu = 1$ as the baseline.

The total weight is:
\begin{equation}
\boxed{
W = 9\wa + \wf + 2\wx + 2\wy + \wix + \wiy + \wdx + \wdy
    + 2\wh + n_s\,\ws + n_u.
}
\label{eq:W}
\end{equation}

The per-byte probability for a byte in class $c$ is:
\begin{equation}
p_c = \frac{w_c}{W}.
\label{eq:pc}
\end{equation}

The minimum byte probability (for unsafe bytes, and for the offset byte) is:
\begin{equation}
\pmin = \frac{\wu}{W} = \frac{1}{W}.
\label{eq:pmin}
\end{equation}

\subsection{Bare core probability per variant}

The 8-byte replicator pattern consists of 6~core bytes, 1~branch
opcode (any of 8, all at probability $\wa/W$), and 1~offset byte
(probability $1/W$).  For each variant at a specific position and
specific rotation, the bare probability is the product of 8 per-byte
probabilities.

From the variant--byte table, each variant uses specific classes.
Note that \hex{00} has probability $\wa/W$ (not a separate class),
and each branch opcode also has probability $\wa/W$.

\paragraph{Standard variants.}
These contain two \hex{00} bytes plus one branch, giving three
factors of $\wa/W$ (i.e.\ $\wa^3/W^3$ from two zeros and one
branch):
\begin{align}
P(\text{DEX bare}) &= \frac{\wx^2 \cdot \wa^3 \cdot \wf \cdot \wdx \cdot \wu}{W^8},
\label{eq:pdex} \\
P(\text{INX bare}) &= \frac{\wx^2 \cdot \wa^3 \cdot \wf \cdot \wix \cdot \wu}{W^8},
\label{eq:pinx} \\
P(\text{DEY bare}) &= \frac{\wy^2 \cdot \wa^3 \cdot \wf \cdot \wdy \cdot \wu}{W^8},
\label{eq:pdey} \\
P(\text{INY bare}) &= \frac{\wy^2 \cdot \wa^3 \cdot \wf \cdot \wiy \cdot \wu}{W^8}.
\label{eq:piny}
\end{align}

\paragraph{Shifted variants.}
These contain one \hex{00} plus one branch ($\wa^2/W^2$), and two
bytes from the shifted-address class ($\wh^2/W^2$):
\begin{align}
P(\text{INX3FF bare}) &= \frac{\wx^2 \cdot \wa^2 \cdot \wix \cdot \wh^2 \cdot \wu}{W^8},
\label{eq:pinx3ff} \\
P(\text{INY3FF bare}) &= \frac{\wy^2 \cdot \wa^2 \cdot \wiy \cdot \wh^2 \cdot \wu}{W^8}.
\label{eq:piny3ff}
\end{align}

\subsection{Incorporating rotations and branch choices}

Each 6-byte core can start at any of 3~instruction boundaries
(rotations).  Each rotation pairs with any of the 8~branch opcodes.
Since all rotations have the same byte content and all 8 branch
opcodes have the same probability $\wa/W$, the multiplicity factor
is $3 \times 8 = 24$.

Summing over all 6 variants, 3 rotations, and 8 branch choices:
\begin{equation}
\boxed{
P_{\text{cell}} = \frac{24\,\wu}{W^8\,(1-S)^7}
  \left[
    \wx^2\,\wa^3\,\wf\,(\wdx + \wix)
  + \wy^2\,\wa^3\,\wf\,(\wdy + \wiy)
  + \wx^2\,\wa^2\,\wix\,\wh^2
  + \wy^2\,\wa^2\,\wiy\,\wh^2
  \right],
}
\label{eq:pcell}
\end{equation}
where $S$ is the survival-weighted safe insert probability
(Section~\ref{sec:inserts}).

\paragraph{Structure of the sum.}
The four terms correspond to:
\begin{enumerate}
\item Standard X-family (DEX + INX): $\wx^2\,\wa^3\,\wf\,(\wdx + \wix)$
\item Standard Y-family (DEY + INY): $\wy^2\,\wa^3\,\wf\,(\wdy + \wiy)$
\item Shifted X (INX3FF): $\wx^2\,\wa^2\,\wix\,\wh^2$
\item Shifted Y (INY3FF): $\wy^2\,\wa^2\,\wiy\,\wh^2$
\end{enumerate}

Factoring:
\begin{equation}
P_{\text{cell}} = \frac{24\,\wu}{W^8\,(1-S)^7}
  \left[
    \wa^2 \left(
      \wa\,\wf\,\bigl[\wx^2(\wdx{+}\wix) + \wy^2(\wdy{+}\wiy)\bigr]
    + \wh^2\,\bigl[\wx^2\,\wix + \wy^2\,\wiy\bigr]
    \right)
  \right].
\label{eq:pcell_factored}
\end{equation}

%----------------------------------------------------------------------
\section{Insert Model}
\label{sec:inserts}

\subsection{Survival-weighted insert probability}

Not all safe bytes are equally likely to preserve viability when
inserted into a loop.  Let $s_c$ denote the \emph{survival weight}
of class $c$ (probability that an insert from class $c$ does not
disrupt the replicator, given execution context).  The effective
safe insert probability per position is:
\begin{equation}
S = \sum_{c \in \text{safe}} n_c \, s_c \, \frac{w_c}{W}
  = \frac{n_s \, s_s \, \ws}{W},
\label{eq:S}
\end{equation}
where the last equality holds when all safe bytes share a single
class with uniform survival weight $s_s$.
In the conservative case ($s_s = 1$), $S = n_s\,\ws / W$.

\subsection{$k$-insert probability}

A replicator with $k$ safe insert bytes has total length $8 + k$.
The $k$ inserts distribute among 7~gaps in the circular loop
(6~core + 1~branch).  The combinatorial factor is:
\begin{equation}
  \binom{k+6}{k}
  \label{eq:nplace}
\end{equation}
(stars-and-bars: $k$ identical slots among 7~gaps).
Each insert byte contributes probability $S$ (survival-weighted).
Thus:
\begin{equation}
P(k\text{-insert}) = \binom{k+6}{k} \cdot P_{\text{bare}} \cdot S^k.
\label{eq:pk}
\end{equation}

\subsection{Total probability at one position}

Summing over $k = 0, 1, 2, \ldots$:
\begin{equation}
P_{\text{pos}} = P_{\text{bare}} \sum_{k=0}^{\infty}
  \binom{k+6}{k} S^k
  = \frac{P_{\text{bare}}}{(1 - S)^7},
\label{eq:ppos}
\end{equation}
using the negative binomial series
$\sum_{k=0}^\infty \binom{k+n-1}{k} x^k = (1-x)^{-n}$ with $n = 7$.

\subsection{Mining time}

The expected number of random boards to mine one replicator on a
$B \times B$ board ($M = 1024$ bytes per cell) is:
\begin{equation}
\boxed{
T = \frac{1}{\rho \cdot P_{\text{cell}}},
}
\label{eq:T}
\end{equation}
where $\rho = B^2 \cdot M$ is the total number of starting positions
on the board ($\rho = 65{,}536 \times 1024 \approx 6.7 \times 10^7$
for $B = 256$).

%----------------------------------------------------------------------
\section{Mean Insert Count}
\label{sec:mean}

Given that a replicator is found, the conditional distribution of
insert count $k$ is:
\begin{equation}
P(k \mid \text{rep}) = \binom{k+6}{k}\,S^k\,(1-S)^7.
\label{eq:pk_cond}
\end{equation}

This is a \emph{negative binomial} distribution with parameters
$r = 7$ and $p = 1 - S$.  The mean is:
\begin{equation}
\boxed{
E[k] = \frac{7\,S}{1 - S}.
}
\label{eq:ek}
\end{equation}

The mean organism length is:
\begin{equation}
\bar{L} = 8 + E[k] \cdot \ell,
\label{eq:Lbar}
\end{equation}
where $\ell$ is the mean number of bytes per insert event
($\ell = 1$ for single-byte safe inserts).

\begin{center}
\begin{tabular}{@{}cc@{}}
\toprule
$S$ & $E[k]$ \\
\midrule
0.01 & 0.07 \\
0.05 & 0.37 \\
0.10 & 0.78 \\
0.20 & 1.75 \\
0.30 & 3.00 \\
0.50 & 7.00 \\
\bottomrule
\end{tabular}
\end{center}

%----------------------------------------------------------------------
\section{Entropy}
\label{sec:entropy}

The Shannon entropy of the soup distribution is:
\begin{equation}
\boxed{
H = \log_2 W - \frac{1}{W}\sum_c n_c\,w_c\,\log_2 w_c.
}
\label{eq:entropy}
\end{equation}

To derive this, note that the per-byte probability for class $c$
is $p_c = w_c/W$, so:
\begin{align}
H &= -\sum_{b=0}^{255} p(b)\,\log_2 p(b)
   = -\sum_c n_c \cdot \frac{w_c}{W}\,\log_2\!\frac{w_c}{W}
  \nonumber \\
  &= -\sum_c \frac{n_c\,w_c}{W}\,(\log_2 w_c - \log_2 W)
  \nonumber \\
  &= \log_2 W \cdot \underbrace{\sum_c \frac{n_c\,w_c}{W}}_{=1}
     - \frac{1}{W}\sum_c n_c\,w_c\,\log_2 w_c.
\label{eq:entropy_deriv}
\end{align}

Maximum entropy is $H_{\max} = \log_2 256 = 8$ bits (all $w_c = 1$,
uniform distribution).  A heavily biased soup has $H \ll 8$.

%----------------------------------------------------------------------
\section{Pareto Analysis}
\label{sec:pareto}

\subsection{The tradeoff triangle}

The three quantities of interest are:
\begin{enumerate}
\item \textbf{Mining time} $T = 1/(\rho \cdot P_{\text{cell}})$:
      lower is better.
\item \textbf{Insert diversity} $E[k]$: higher means richer
      organisms with more ``junk DNA.''
\item \textbf{Entropy} $H$: higher means the soup looks more random
      and produces more diverse non-replicator cells.
\end{enumerate}

\subsection{Key tension: the offset byte floor}

The offset byte probability is $\pmin = 1/W$.  Increasing any class
weight $w_c$ raises $W$, which \emph{lowers} $\pmin$.  Since every
replicator includes one offset byte, the mining probability contains
a factor of $1/W$.  Combined with the $1/W^8$ from the 8-byte product,
the total scaling is $1/W^8$ for the bare core.

Concentrating the soup (raising core weights) simultaneously:
\begin{itemize}
\item \emph{Increases} the numerator (products of core class weights),
\item \emph{Increases} the denominator ($W^8$),
\end{itemize}
creating a self-regulating balance.  The offset byte acts as a
\emph{tax} on enrichment that prevents infinite concentration from
being optimal.

\subsection{Symmetric simplification}

Setting $\wx = \wy$, $\wix = \wiy = \wdx = \wdy$, and treating
the shifted variants as negligible ($\wh = 1$, contributing
little to the sum), the formula reduces to \textbf{5 free parameters}:
$\wa, \wf, w_{xy} \equiv \wx = \wy,
 w_{\mathrm{inc}} \equiv \wix = \wiy = \wdx = \wdy, \ws$.

In this regime:
\begin{align}
W &= 9\wa + \wf + 4\,w_{xy} + 4\,w_{\mathrm{inc}} + 2
   + n_s\,\ws + n_u, \\
P_{\text{cell}} &\propto
  \frac{w_{xy}^2 \cdot \wa^3 \cdot \wf \cdot 2\,w_{\mathrm{inc}}}
       {W^8\,(1-S)^7}. \label{eq:pcell_sym}
\end{align}

\subsection{Optimal core allocation}

The exponent structure of Eq.~\eqref{eq:pcell_sym} shows that the
mining probability is a monomial in the class weights:
$\wa^3 \cdot w_{xy}^2 \cdot \wf^1 \cdot w_{\mathrm{inc}}^1$
divided by $W^8$.  By the AM--GM inequality (or Lagrange multipliers
on the constraint $W = \text{const}$), the optimal allocation is
proportional to the exponents:
\begin{equation}
9\wa : 4w_{xy} : \wf : 4w_{\mathrm{inc}} \;\propto\; 3 : 2 : 1 : 1.
\end{equation}
That is, allocate weight in proportion to the exponent with which
each class contributes, divided by its multiplicity in $W$.

\subsection{Key insight: class value scales with variant coverage}

Byte \hex{04} (class $\wf$, needed by 4 standard variants) is
more valuable per byte than \hex{CA} (class $\wdx$, needed by
1 variant only).  This is because $\wf$ appears in 4 terms of
the $P_{\text{cell}}$ sum (all standard variants), while $\wdx$
appears in only 1 term (DEX).  The power-set structure of the
classes directly encodes this differential value.

%----------------------------------------------------------------------
\section{Discussion}
\label{sec:discussion}

The power-set soup class framework reveals several insights:

\begin{enumerate}
\item \textbf{The offset byte tax.}
  The offset byte has probability $\pmin = 1/W$.  As core weights
  increase, $W$ grows and $\pmin$ shrinks.  This creates a
  self-regulating mechanism: concentrating the soup \emph{lowers}
  the offset probability, counteracting the benefit of enriching
  core bytes.  The net effect is that $T$ has a finite minimum
  as a function of the weights.

\item \textbf{Power-set class values.}
  The proper partition groups bytes by the \emph{subset of variants
  that need them}.  This reveals a hierarchy of value:
  $\wa$ (all variants) $>$ $\wf$ (4 variants) $>$ $\wx/\wy$
  (3 variants each) $>$ $\wix/\wiy$ (2 variants) $>$
  $\wdx/\wdy$ (1 variant) $>$ $\wh$ (2 shifted variants).

\item \textbf{Symmetry breaking is wasteful.}
  The optimal allocation has $\wx = \wy$ and
  $\wix = \wiy = \wdx = \wdy$.
  Breaking X/Y symmetry concentrates weight on fewer variant terms,
  reducing total probability.

\item \textbf{Free inserts.}
  When safe insert weight is increased, the $(1-S)^{-7}$
  amplification from the negative binomial compensates for the
  increase in $W$.  Inserts enrich the organism population at
  modest mining cost, especially when $S \ll 1$.

\item \textbf{Entropy is expensive.}
  The dominant factor in mining difficulty is the unsafe count
  $n_u$, which contributes $n_u$ to $W$ at the baseline weight
  of 1.  The mining time scales as $W^8$, so each additional
  unit of $W$ from unsafe bytes is multiplicatively costly.
\end{enumerate}

%----------------------------------------------------------------------
\appendix
\section{Complete Byte Classification}
\label{app:bytes}

Table~\ref{tab:bytes} classifies all 256 byte values.  The
\textbf{Class} column uses abbreviations:
\textbf{a} = all variants + branches ($\wa$),
\textbf{f} = standard page ($\wf$),
\textbf{x} = X load/store ($\wx$),
\textbf{y} = Y load/store ($\wy$),
\textbf{ix} = X increment ($\wix$),
\textbf{iy} = Y increment ($\wiy$),
\textbf{dx} = X decrement ($\wdx$),
\textbf{dy} = Y decrement ($\wdy$),
\textbf{h} = shifted address ($\wh$),
\textbf{s} = safe inserts ($\ws$),
\textbf{u} = unsafe ($\wu$).
\textbf{Len} = instruction length (1, 2, or 3 bytes; 0 = JAM).
\textbf{Mnemonic} shows the official or common undocumented name.

{\small
\begin{longtable}{@{}ccclp{4.8cm}@{}}
\caption{Byte classification for all 256 values.}
\label{tab:bytes} \\
\toprule
\textbf{Hex} & \textbf{Class} & \textbf{Len} & \textbf{Mnemonic} &
\textbf{Notes} \\
\midrule
\endfirsthead
\toprule
\textbf{Hex} & \textbf{Class} & \textbf{Len} & \textbf{Mnemonic} &
\textbf{Notes} \\
\midrule
\endhead
\midrule
\multicolumn{5}{r}{\emph{continued}} \\
\endfoot
\bottomrule
\endlastfoot
% Row 0x
00 & a  & 1 & BRK & Core (all variants) \\
01 & u  & 2 & ORA (zp,X) & \\
02 & u  & 0 & JAM & CPU halt \\
03 & h  & 2 & SLO (zp,X) & Core (shift variants) \\
04 & f  & 2 & NOP zp & Core (std variants) \\
05 & u  & 2 & ORA zp & \\
06 & u  & 2 & ASL zp & \\
07 & u  & 2 & SLO zp & \\
08 & s  & 1 & PHP & Safe: push P \\
09 & u  & 2 & ORA \#imm & \\
0A & u  & 1 & ASL A & Clobbers A, flags \\
0B & u  & 2 & ANC \#imm & \\
0C & u  & 3 & NOP abs & \\
0D & u  & 3 & ORA abs & \\
0E & u  & 3 & ASL abs & \\
0F & u  & 3 & SLO abs & \\
% Row 1x
10 & a  & 2 & BPL & Branch opcode \\
11 & u  & 2 & ORA (zp),Y & \\
12 & u  & 0 & JAM & CPU halt \\
13 & u  & 2 & SLO (zp),Y & \\
14 & u  & 2 & NOP zp,X & \\
15 & u  & 2 & ORA zp,X & \\
16 & u  & 2 & ASL zp,X & \\
17 & u  & 2 & SLO zp,X & \\
18 & u  & 1 & CLC & Clobbers C flag \\
19 & u  & 3 & ORA abs,Y & \\
1A & s  & 1 & NOP & Safe: undoc NOP \\
1B & u  & 3 & SLO abs,Y & \\
1C & u  & 3 & NOP abs,X & \\
1D & u  & 3 & ORA abs,X & \\
1E & u  & 3 & ASL abs,X & \\
1F & u  & 3 & SLO abs,X & \\
% Row 2x
20 & u  & 3 & JSR abs & \\
21 & u  & 2 & AND (zp,X) & \\
22 & u  & 0 & JAM & CPU halt \\
23 & u  & 2 & RLA (zp,X) & \\
24 & u  & 2 & BIT zp & \\
25 & u  & 2 & AND zp & \\
26 & u  & 2 & ROL zp & \\
27 & u  & 2 & RLA zp & \\
28 & u  & 1 & PLP & Clobbers all flags \\
29 & u  & 2 & AND \#imm & \\
2A & u  & 1 & ROL A & Clobbers A, flags \\
2B & u  & 2 & ANC \#imm & \\
2C & u  & 3 & BIT abs & \\
2D & u  & 3 & AND abs & \\
2E & u  & 3 & ROL abs & \\
2F & u  & 3 & RLA abs & \\
% Row 3x
30 & a  & 2 & BMI & Branch opcode \\
31 & u  & 2 & AND (zp),Y & \\
32 & u  & 0 & JAM & CPU halt \\
33 & u  & 2 & RLA (zp),Y & \\
34 & u  & 2 & NOP zp,X & \\
35 & u  & 2 & AND zp,X & \\
36 & u  & 2 & ROL zp,X & \\
37 & u  & 2 & RLA zp,X & \\
38 & u  & 1 & SEC & Clobbers C flag \\
39 & u  & 3 & AND abs,Y & \\
3A & s  & 1 & NOP & Safe: undoc NOP \\
3B & u  & 3 & RLA abs,Y & \\
3C & u  & 3 & NOP abs,X & \\
3D & u  & 3 & AND abs,X & \\
3E & u  & 3 & ROL abs,X & \\
3F & u  & 3 & RLA abs,X & \\
% Row 4x
40 & u  & 1 & RTI & Control flow \\
41 & u  & 2 & EOR (zp,X) & \\
42 & u  & 0 & JAM & CPU halt \\
43 & u  & 2 & SRE (zp,X) & \\
44 & u  & 2 & NOP zp & \\
45 & u  & 2 & EOR zp & \\
46 & u  & 2 & LSR zp & \\
47 & u  & 2 & SRE zp & \\
48 & s  & 1 & PHA & Safe: push A \\
49 & u  & 2 & EOR \#imm & \\
4A & u  & 1 & LSR A & Clobbers A, flags \\
4B & u  & 2 & ALR \#imm & \\
4C & u  & 3 & JMP abs & Control flow \\
4D & u  & 3 & EOR abs & \\
4E & u  & 3 & LSR abs & \\
4F & u  & 3 & SRE abs & \\
% Row 5x
50 & a  & 2 & BVC & Branch opcode \\
51 & u  & 2 & EOR (zp),Y & \\
52 & u  & 0 & JAM & CPU halt \\
53 & u  & 2 & SRE (zp),Y & \\
54 & u  & 2 & NOP zp,X & \\
55 & u  & 2 & EOR zp,X & \\
56 & u  & 2 & LSR zp,X & \\
57 & u  & 2 & SRE zp,X & \\
58 & u  & 1 & CLI & Clobbers I flag (safe for branches, but conservatively excluded) \\
59 & u  & 3 & EOR abs,Y & \\
5A & s  & 1 & NOP & Safe: undoc NOP \\
5B & u  & 3 & SRE abs,Y & \\
5C & u  & 3 & NOP abs,X & \\
5D & u  & 3 & EOR abs,X & \\
5E & u  & 3 & LSR abs,X & \\
5F & u  & 3 & SRE abs,X & \\
% Row 6x
60 & u  & 1 & RTS & Control flow \\
61 & u  & 2 & ADC (zp,X) & \\
62 & u  & 0 & JAM & CPU halt \\
63 & u  & 2 & RRA (zp,X) & \\
64 & u  & 2 & NOP zp & \\
65 & u  & 2 & ADC zp & \\
66 & u  & 2 & ROR zp & \\
67 & u  & 2 & RRA zp & \\
68 & u  & 1 & PLA & Clobbers A, flags \\
69 & u  & 2 & ADC \#imm & \\
6A & u  & 1 & ROR A & Clobbers A, flags \\
6B & u  & 2 & ARR \#imm & \\
6C & u  & 3 & JMP (abs) & Control flow \\
6D & u  & 3 & ADC abs & \\
6E & u  & 3 & ROR abs & \\
6F & u  & 3 & RRA abs & \\
% Row 7x
70 & a  & 2 & BVS & Branch opcode \\
71 & u  & 2 & ADC (zp),Y & \\
72 & u  & 0 & JAM & CPU halt \\
73 & u  & 2 & RRA (zp),Y & \\
74 & u  & 2 & NOP zp,X & \\
75 & u  & 2 & ADC zp,X & \\
76 & u  & 2 & ROR zp,X & \\
77 & u  & 2 & RRA zp,X & \\
78 & u  & 1 & SEI & Clobbers I flag (conservatively excluded) \\
79 & u  & 3 & ADC abs,Y & \\
7A & s  & 1 & NOP & Safe: undoc NOP \\
7B & u  & 3 & RRA abs,Y & \\
7C & u  & 3 & NOP abs,X & \\
7D & u  & 3 & ADC abs,X & \\
7E & u  & 3 & ROR abs,X & \\
7F & u  & 3 & RRA abs,X & \\
% Row 8x
80 & u  & 2 & NOP \#imm & Undoc 2-byte NOP \\
81 & u  & 2 & STA (zp,X) & \\
82 & u  & 2 & NOP \#imm & Undoc 2-byte NOP \\
83 & u  & 2 & SAX (zp,X) & \\
84 & u  & 2 & STY zp & \\
85 & u  & 2 & STA zp & ZP write \\
86 & u  & 2 & STX zp & ZP write \\
87 & u  & 2 & SAX zp & ZP write \\
88 & dy & 1 & DEY & Core (DEY variant) \\
89 & u  & 2 & NOP \#imm & Undoc 2-byte NOP \\
8A & u  & 1 & TXA & Clobbers A, flags \\
8B & u  & 2 & ANE \#imm & \\
8C & u  & 3 & STY abs & \\
8D & u  & 3 & STA abs & \\
8E & u  & 3 & STX abs & \\
8F & u  & 3 & SAX abs & \\
% Row 9x
90 & a  & 2 & BCC & Branch opcode \\
91 & u  & 2 & STA (zp),Y & \\
92 & u  & 0 & JAM & CPU halt \\
93 & u  & 2 & SHA (zp),Y & \\
94 & u  & 2 & STY zp,X & \\
95 & u  & 2 & STA zp,X & ZP write \\
96 & u  & 2 & STX zp,Y & ZP write \\
97 & u  & 2 & SAX zp,Y & ZP write \\
98 & u  & 1 & TYA & Clobbers A, flags \\
99 & y  & 3 & STA abs,Y & Core (Y variants) \\
9A & u  & 1 & TXS & Modifies only S (no flags), but not a NOP-like operation \\
9B & u  & 3 & TAS abs,Y & \\
9C & u  & 3 & SHY abs,X & \\
9D & x  & 3 & STA abs,X & Core (X variants) \\
9E & u  & 3 & SHX abs,Y & \\
9F & u  & 3 & SHA abs,Y & \\
% Row Ax
A0 & u  & 2 & LDY \#imm & Clobbers Y, flags \\
A1 & u  & 2 & LDA (zp,X) & \\
A2 & u  & 2 & LDX \#imm & Clobbers X, flags \\
A3 & u  & 2 & LAX (zp,X) & \\
A4 & u  & 2 & LDY zp & \\
A5 & u  & 2 & LDA zp & \\
A6 & u  & 2 & LDX zp & \\
A7 & u  & 2 & LAX zp & \\
A8 & u  & 1 & TAY & Clobbers Y, N, Z \\
A9 & u  & 2 & LDA \#imm & \\
AA & u  & 1 & TAX & Clobbers X, N, Z \\
AB & u  & 2 & LAX \#imm & \\
AC & u  & 3 & LDY abs & \\
AD & u  & 3 & LDA abs & \\
AE & u  & 3 & LDX abs & \\
AF & u  & 3 & LAX abs & \\
% Row Bx
B0 & a  & 2 & BCS & Branch opcode \\
B1 & u  & 2 & LDA (zp),Y & \\
B2 & u  & 0 & JAM & CPU halt \\
B3 & u  & 2 & LAX (zp),Y & \\
B4 & u  & 2 & LDY zp,X & \\
B5 & x  & 2 & LDA zp,X & Core (X variants) \\
B6 & u  & 2 & LDX zp,Y & \\
B7 & y  & 2 & LAX zp,Y & Core (Y variants) \\
B8 & u  & 1 & CLV & Clobbers V flag \\
B9 & u  & 3 & LDA abs,Y & \\
BA & u  & 1 & TSX & Clobbers X, N, Z \\
BB & u  & 3 & LAS abs,Y & \\
BC & u  & 3 & LDY abs,X & \\
BD & u  & 3 & LDA abs,X & \\
BE & u  & 3 & LDX abs,Y & \\
BF & u  & 3 & LAX abs,Y & \\
% Row Cx
C0 & u  & 2 & CPY \#imm & \\
C1 & u  & 2 & CMP (zp,X) & \\
C2 & u  & 2 & NOP \#imm & Undoc 2-byte NOP \\
C3 & u  & 2 & DCP (zp,X) & \\
C4 & u  & 2 & CPY zp & \\
C5 & u  & 2 & CMP zp & \\
C6 & u  & 2 & DEC zp & ZP write \\
C7 & u  & 2 & DCP zp & ZP write \\
C8 & iy & 1 & INY & Core (INY, INY3FF) \\
C9 & u  & 2 & CMP \#imm & \\
CA & dx & 1 & DEX & Core (DEX variant) \\
CB & u  & 2 & AXS \#imm & \\
CC & u  & 3 & CPY abs & \\
CD & u  & 3 & CMP abs & \\
CE & u  & 3 & DEC abs & \\
CF & u  & 3 & DCP abs & \\
% Row Dx
D0 & a  & 2 & BNE & Branch opcode \\
D1 & u  & 2 & CMP (zp),Y & \\
D2 & u  & 0 & JAM & CPU halt \\
D3 & u  & 2 & DCP (zp),Y & \\
D4 & u  & 2 & NOP zp,X & \\
D5 & u  & 2 & CMP zp,X & \\
D6 & u  & 2 & DEC zp,X & ZP write \\
D7 & u  & 2 & DCP zp,X & \\
D8 & u  & 1 & CLD & Clobbers D flag (conservatively excluded) \\
D9 & u  & 3 & CMP abs,Y & \\
DA & s  & 1 & NOP & Safe: undoc NOP \\
DB & u  & 3 & DCP abs,Y & \\
DC & u  & 3 & NOP abs,X & \\
DD & u  & 3 & CMP abs,X & \\
DE & u  & 3 & DEC abs,X & \\
DF & u  & 3 & DCP abs,X & \\
% Row Ex
E0 & u  & 2 & CPX \#imm & \\
E1 & u  & 2 & SBC (zp,X) & \\
E2 & u  & 2 & NOP \#imm & Undoc 2-byte NOP \\
E3 & u  & 2 & ISC (zp,X) & \\
E4 & u  & 2 & CPX zp & \\
E5 & u  & 2 & SBC zp & \\
E6 & u  & 2 & INC zp & ZP write \\
E7 & u  & 2 & ISC zp & ZP write \\
E8 & ix & 1 & INX & Core (INX, INX3FF) \\
E9 & u  & 2 & SBC \#imm & \\
EA & s  & 1 & NOP & Safe: official NOP \\
EB & u  & 2 & SBC \#imm & Undoc duplicate \\
EC & u  & 3 & CPX abs & \\
ED & u  & 3 & SBC abs & \\
EE & u  & 3 & INC abs & \\
EF & u  & 3 & ISC abs & \\
% Row Fx
F0 & a  & 2 & BEQ & Branch opcode \\
F1 & u  & 2 & SBC (zp),Y & \\
F2 & u  & 0 & JAM & CPU halt \\
F3 & u  & 2 & ISC (zp),Y & \\
F4 & u  & 2 & NOP zp,X & \\
F5 & u  & 2 & SBC zp,X & \\
F6 & u  & 2 & INC zp,X & ZP write \\
F7 & u  & 2 & ISC zp,X & \\
F8 & u  & 1 & SED & Clobbers D flag (conservatively excluded) \\
F9 & u  & 3 & SBC abs,Y & \\
FA & s  & 1 & NOP & Safe: undoc NOP \\
FB & u  & 3 & ISC abs,Y & \\
FC & u  & 3 & NOP abs,X & \\
FD & u  & 3 & SBC abs,X & \\
FE & u  & 3 & INC abs,X & \\
FF & h  & 3 & ISC abs,X & Core (shift variants) \\
\end{longtable}
}

\paragraph{Summary.}
Class $a$ (all + branches): 9 values (1 core + 8 branch).
Classes $x, y$: 2 values each (load/store pairs).
Classes $ix, iy, dx, dy$: 1 value each (increment/decrement).
Class $f$ (standard page): 1 value.
Class $h$ (shifted address): 2 values.
Safe inserts ($s$): 9 values (conservative).
Unsafe ($u$): 227 values.
Total: $9 + 2 + 2 + 1 + 1 + 1 + 1 + 1 + 2 + 9 + 227 = 256$.\ \checkmark

\paragraph{A note on TXS (\hex{9A}).}
\opcode{TXS} modifies only S with no flag side effects, so it is
arguably safe.  We conservatively exclude it because transferring a
loop counter to the stack pointer could cause erratic stack behavior
in the presence of PHA/PHP inserts.  Including it would increase
$|\Csafe|$ to 10.

\paragraph{A note on SEI/CLI/SED/CLD (\hex{78}/\hex{58}/\hex{F8}/\hex{D8}).}
These modify only the I or D flags, which are not tested by any
conditional branch.  Including them would not affect branch behavior
and could be considered safe.  We exclude them from the conservative
set; a relaxed analysis could add 4 more bytes to $\Csafe$ (total 13).

\end{document}
